\section{Worked realizations in the common operator language}

\subsection{Signal polyphase coordinates}

Let $x=(x_n)_{n\in\Z}\in c_{00}(\Z;R)$ and let the seed be $\delta_0$
under the bilateral shifts $(S^r\delta_0)_n=\delta_{n,r}$. This makes the
algebraic span and every displayed sum finite. The $N$ packets are
\[
x_j[m]=x[Nm+j].
\]
Synthesis interleaves these sequences. A linear filter with a finite impulse
response that is invariant under shifts becomes an $N\times N$ matrix of
transfer blocks. Refining from $N$ to $M$ with $N\mid M$ splits each coarse
phase into its descendants without changing the reconstructed signal. The
same formulas extend to standard sequence completions when the synthesis
series converges. This is the usual polyphase viewpoint from multirate filter
banks \cite{Vaidyanathan1993}.

\subsection{Sparsity of transfer blocks}

A banded operator need not remain banded after residue reordering, but its
nonzero couplings become a structured packet graph. With the convention
$(S_rx)_n=x_{n-r}$, the only active transfer is
\[
P_{N,j+r}S_rP_{N,j},
\]
with indices taken modulo $N$. More general finite stencils produce a finite set of active phase differences.

\begin{figure}[htbp]
\centering
\includegraphics[width=0.96\textwidth]{10_transfer_block_sparsity.pdf}
\caption{Norms of packet blocks for a finite stencil. The original operator is exactly reconstructed from a sparse transfer pattern.}
\label{fig:block-sparsity}
\end{figure}

\subsection{A short realization in Smith chart coordinates}\label{sec:smith-realization}

\begin{interpretation}
The Smith chart used here is a coordinate diagram for impedance and reflection. It is different from the Smith normal form used earlier to study integral congruences. The connection within this paper is the way a change of coordinates makes an action easier to describe. Here a rotation in reflection coordinates becomes a fractional linear map in impedance coordinates. The sampled phase orbit is a geometric example and does not supply another linear projector decomposition.
\end{interpretation}

Let $Z$ and $Z_0\ne0$ be complex impedances, put $z=Z/Z_0$, and assume $z\ne-1$. Write
\[
\Gamma=\frac{z-1}{z+1}.
\]
On the open reflection disk $|\Gamma|<1$, multiplication by any phase $\omega_N^k$ stays in the disk and has a finite inverse-chart image. The phase action is
\[
\Gamma\mapsto\omega_N^k\Gamma.
\]
Outside the open disk, the same chart formula requires $1-\omega_N^k\Gamma\ne0$. Pulling the phase action back on this domain gives
\[
T_{N,k}(z)
=
\frac{(1+\omega_N^k)z+(1-\omega_N^k)}
{(1-\omega_N^k)z+(1+\omega_N^k)}.
\]
For $N=2$ and $k=1$,
\phantomsection
\[
T_{2,1}(z)=\frac1z,
\qquad z\ne0,-1.
\label{eq:smith-inversion}
\]
Thus the nontrivial action at the first phase level exchanges normalized
impedance and admittance. Calling these values ``packets'' is geometric
terminology for the sampled orbit. It does not assert that the two values
form linear idempotent ranges. Higher refinements at powers of $2$ sample the
same reflection circle at progressively finer electrical phases. Compare the
standard Smith chart conventions in \cite{pozar2012}.

\begin{figure}[htbp]
\centering
\includegraphics[width=0.96\textwidth]{11_smith_polyphase_orbit.pdf}
\caption{Nested phase orbits at powers of $2$ on the Smith reflection disk. The full treatment of the Smith chart is deferred to the signal paper.}
\label{fig:smith-orbit}
\end{figure}

The symmetric weight
\[
w_{\mathrm{Smith}}(\Gamma)
=
\frac1{|1-\Gamma|^2}
+
\frac1{|1+\Gamma|^2}
\]
indicates where small errors in reflection coordinates produce large impedance or admittance errors.

\begin{figure}[htbp]
\centering
\includegraphics[width=0.96\textwidth]{12_smith_adaptive_weight.pdf}
\caption{A symmetric refinement weight concentrated near the boundary points for open and short circuits.}
\label{fig:smith-weight}
\end{figure}


\section{Finite computational certificates}
The following results and tables record finite checks of the implementation. They do not certify the general lattice or analytic statements in Lean. The accompanying verification files record the execution evidence and tested ranges.


All numerical work in this paper is finite and reproducible. It tests the
Python implementation of identities proved above or checked against
certificates. It does not replace those arguments or the Lean kernel checks.
The recorded implementation run covers 48 tests, including contracts for invalid input, refinement
transitivity, aggregation of transfer blocks, Fourier/DFT agreement, Smith
fibre counts by exhaustive enumeration, and consistency of the stored JSON
and CSV records.

\begin{table}[htbp]
\centering
\caption{Representative numerical and symbolic checks.}
\begin{tabularx}{\textwidth}{L{0.54\textwidth}Y}
\toprule
Check & Result \\
\midrule
Packet synthesis residuals for moduli $2,3,4,6,8,12$ & $0$ \\
Nested refinement residuals for $2\to4$, $2\to8$, $3\to6$, $3\to12$, $4\to8$, $4\to12$, $6\to12$ & $0$ \\
CRT residuals for dyadic and odd moduli for $4\times3$, $8\times3$, $4\times5$, $8\times5$ & $0$ \\
Maximum finite Fourier projector residual & $6.60\times10^{-15}$ \\
Weighted inverse residual modulo $2^{10}$ & $0$ \\
Weighted inverse residual modulo $3^8$ & $0$ \\
Weighted inverse residual modulo $5^7$ & $0$ \\
Difference between filtered and direct modular inverses & $0$ \\
Transfer zeta coefficients versus direct walk enumeration through weight $30$ & exact match \\
Dyadic sign projector resultant weight & $v_2(2)=1$ \\
Sign projector resultant weight at the odd prime $3$ & $v_3(2)=0$ \\
Power defect telescoping identity residual & at most $6.94\times10^{-18}$ \\
Resolvent defect identity residual & $1.09\times10^{-17}$ \\
Smith $N=2$ phase versus $z\mapsto1/z$ & residual below $1.4\times10^{-16}$ \\
\bottomrule
\end{tabularx}
\end{table}

\begin{certificate}[Weighted inverse]
For the tested diagonal-plus-transfer matrices, the filtered inverse agrees exactly with direct modular inversion at all tested precisions:
\[
2^{10},\qquad3^8,\qquad5^7.
\]
The maximum modular residual is zero in every case.
\end{certificate}

\begin{certificate}[Transfer zeta]
The rational function
\[
\frac{-t^3}{t^9-t^7-t^6-t^5+2t^4+t^3+t^2-1}
\]
was expanded through degree $30$. Every coefficient agrees with direct weighted-walk enumeration.
\end{certificate}

\begin{certificate}[Dyadic separation]
The code verifies exactly that the two rational sign projectors are idempotent, orthogonal, and complementary, that scaling by $2$ makes both integral, and that the same projectors are already integral modulo every tested odd prime.
\end{certificate}


\section{Scope and limits}\label{sec:limits}

The exact statements above have several deliberate boundaries.

First, refinement changes coordinates while preserving reconstruction. If a numerical method discards fine packets, the resulting error must be estimated separately.

Second, the weighted inverse is linear.  A nonlinear map creates higher order interactions among packet coordinates.  Those interactions are not contained in the path formula of \cref{thm:weighted_inverse}.

Third, a resultant is only a denominator bound.  The exact weight depends on the lattice and is given by \cref{thm:gluing_weight}.  The two examples in \cref{ex:loose_resultant,ex:lattice_dependence} show that the same eigenvalue separation can give different integral packet weights.

Fourth, the dyadic sign example is a model rather than a classification of all phenomena at the prime two.  Character averages, trace splittings, quadratic refinements, and deformation projectors must each be checked on their own lattices.

Finally, the defect formulas do not prove that a chosen coarse operator is a good approximation.  They say exactly how a known one step defect propagates.  A useful approximation theorem must also bound that defect in the problem under study.



\appendix
\section{Notation table}

\begin{longtable}{L{0.28\textwidth}L{0.64\textwidth}}
\toprule
Symbol & Meaning \\
\midrule
\endhead
$P_{\ell,\alpha}$ & packet projector at refinement level $\ell$ \\
$I_\ell$ & packet address set at level $\ell$ \\
$\beta\succ\alpha$ & fine packet $\beta$ is a child of coarse packet $\alpha$ \\
$\cA_\ell,\cS_\ell$ & packet analysis and synthesis maps \\
$J_{\ell+1,\ell}$ & exact coordinate refinement map \\
$C_{\ell,\ell+1}$ & exact coordinate regrouping from fine to coarse \\
$T^{(\ell)}_{\beta\alpha}$ & transfer block $P_{\ell,\beta}TP_{\ell,\alpha}$ \\
$P_{N,j}$ & residue or phase selector for address $j\pmod N$ \\
$w(e)$ & positive filtration weight of a transfer edge \\
$W(\gamma)$ & cumulative weight of a transfer word or graph walk \\
$A(t)$ & polynomial weighted adjacency matrix \\
$Z_{s,t}(t)$ & weighted transfer-zeta function from vertex $s$ to vertex $t$ \\
$s_{\pi,\Lambda}(E)$ & separation weight of a rational projector on the chosen lattice, when a clearing exponent exists \\
$\Delta$ & linear refinement defect for one step $CT_f-T_cC$ \\
\bottomrule
\end{longtable}

\section{Algorithmic forms}\label{sec:algorithms}

\subsection{Residue analysis and synthesis}\label{alg:residue}

\begin{lstlisting}
def analyze_residues(x, N):
    require N > 0
    return [x[j::N] for j in range(N)]

def synthesize_residues(packets):
    N = len(packets)
    require packet lengths fit one N-way interleaving
    y = empty(sum(len(p) for p in packets))
    for j, packet in enumerate(packets):
        y[j::N] = packet
    return y
\end{lstlisting}

The verified contract is
$\operatorname{synthesize}(\operatorname{analyze}(x,N))=x$. The Lean
counterpart uses packet coordinates restricted to projector ranges and proves
both inverse laws. The finite Python routine checks the concrete array
realization.

\subsection{Weighted inverse at finite precision}\label{alg:weighted-inverse}

\begin{lstlisting}
Dinv = inverse(D)                       # certified on both sides
edge[e] = Dinv @ H_edge[e]
require Dinv @ H == sum(edge[e] for e in edges)
require edge[e] lies in filtration[weight[e]]
m = certified_lower_bound(weight)       # 1 <= m <= weight[e]
N = (L + m - 1) // m                    # so L <= m*N
S = 0
for n in range(N):
    for word in edge_tuples(edges, n):
        if sum(weight[e] for e in word) < L:
            S += (-1)**n * ordered_product(edge[e] for e in word)
return q_L(S @ Dinv)
\end{lstlisting}

Lean proves that the mapped pruned sum is exactly the mapped finite geometric
inverse. The lower bound $m$ gives $m|\gamma|\le W(\gamma)$, so
$L\le mN$ derives $q_L(A^N)=0$ from the filtration certificate and makes the
returned value an inverse on both sides of $q_L(D+H)$. The Lean lemma
$L\le m\lceil L/m\rceil$ justifies the displayed executable cutoff for
$m>0$. Taking $m=1$ and $N=L$ gives
the unconditional cutoff for positive weights. The Python implementation checks the
DVR matrix specialization against direct modular inversion.

\subsection{Walk construction for transfer zeta}\label{alg:path-zeta}

\begin{lstlisting}
def coefficient(G, start, target, n):
    total = 0
    for length in range(n + 1):
        for word in edge_tuples(G, length):
            if is_walk(word, start, target) and weight(word) == n:
                total += 1
    return total

A[i,j] = sum(t**weight(e) for distinct edges e from i to j)
T = I - A
Q = det(T)
H = adjugate(T)
P = H[start, target]
assert Q.coeff(0) == 1
Z = coefficient_series(G, start, target)
prove Q * Z == P
\end{lstlisting}

Positive weights make the coefficient loop finite because every contributing
word has length at most $n$. Lean proves the enumeration correct, derives its
recurrence from the first edge, proves $(I-A)Z=I$, constructs the
determinant/adjugate inverse, and identifies the two matrices by inverse
uniqueness. Thus the normalized $QZ=P$ identity and its induced recurrence
follow from finite graph data. The Python companion repeats the exact symbolic
construction as an executable check. The Lean proof does not use it as an
oracle.

\section{Proof dependency map}

\begin{figure}[htbp]
\centering
\includegraphics[width=0.96\textwidth]{13_theorem_dependency.pdf}
\caption{The dependency structure of Chapter II. Coordinate coherence, residue/CRT, weighted inversion, automatic finite-graph transfer rationality, and integral separation are distinct branches with explicit inputs.}
\label{fig:dependency}
\end{figure}

The proof order is:
\[
\begin{aligned}
&\text{complete packet axioms}
  \Rightarrow\text{finite-packet uniqueness and coordinate coherence}\\
&\qquad\Rightarrow\text{coordinate naturality and refined transfer blocks},\\
&\text{residue selectors and a finite Fourier table}
  \Rightarrow\text{Fourier formulas},\\
&\text{coprime moduli}\Rightarrow\text{dyadic--odd formulas},\\
&\text{quotient, filtration, and edge-weight certificates}
  \Rightarrow\text{explicit pruned inverse},\\
&\text{finite positive-weight graph}
  \Rightarrow\text{semantic enumeration and first-edge recurrence}\\
&\qquad\Rightarrow\det/\operatorname{adj}\text{ rationality and recurrence},\\
&\text{minimal evaluated B\'ezout data and inverse-unit lattice action}\\
&\qquad\Rightarrow\text{separation bound, upward closure, and precision delay},\\
&\text{concrete }\mathbb Z/p^L\mathbb Z\text{ coordinate fibres}
  \Rightarrow\text{finite Smith product count},\\
&\text{one-step typed defect}
  \Rightarrow\text{powers, polynomial actions, and witnessed resolvents}.
\end{aligned}
\]

The walk branch for transfer zeta is logically independent of operator
inversion. Its rationality conclusion follows from finite graph data without
a supplied $QZ=P$ premise. Likewise, CRT factorization follows directly from
residue selectors and coprimality. It does not require the abstract refinement
theorem, although it supplies a concrete refinement family. General Smith
diagonalization and the construction of polynomial B\'ezout/resultant data
remain inputs to the formalized branches that use those certificates.

